\section*{Chapitre \ref{ch:m1:auctions} --- Enchères}

\begin{solution}{exo:m1:auctions:1}
À 19.95 et 20.00~: $D = 1\,200$, $S = 200$, volume exécutable 200. À 20.05~: $D =
700$, $S = 800$, volume exécutable 700. À 20.10~: $D = 400$, $S = 1\,100$,
volume exécutable 400. Le \omterm{def:m1:auctions:call}{prix d'équilibre} est 20.05, pour 700 actions échangées.
\end{solution}

\begin{solution}{exo:m1:auctions:2}
Excédent $700 - 800 = -100$~: 100 actions restent à vendre à 20.05. Les deux
ordres d'achat à 20.10 et 20.05 sont exécutés intégralement~; l'ordre de vente à 19.95 l'est aussi intégralement
et celui à 20.05 pour 500 de ses 600. Les ordres à 20.00 (achat) et à 20.10
(vente) ne donnent lieu à aucune transaction.
\end{solution}

\begin{solution}{exo:m1:auctions:3}
Sur Nasdaq, librement~: la saisie des ordres MOC prend fin à 15:55. À la New York Stock Exchange
seulement si un déséquilibre réglementaire \emph{vendeur} a été publié à 15:50 pour ce
symbole~: après l'heure limite, les ordres de clôture ne sont acceptés que du côté
qui compense le déséquilibre publié.
\end{solution}

\begin{solution}{exo:m1:auctions:4}
Chaque prix compris entre 50.00 et 50.20 exécute 1\,000 actions, sans excédent~:
les règles 1 à 3 ne tranchent pas. La règle 4 donne le prix de référence, 50.12. Avec un
prix de référence de 50.40, hors de la plage, le candidat le plus proche est 50.20. Deux
parties prêtes à négocier n'importe où dans une plage de 20 cents voient leur prix fixé par la dernière
transaction, ni par l'une ni par l'autre.
\end{solution}

\begin{solution}{exo:m1:auctions:5}
$300\,000/(25\,000 + 40\,000) = 4.6$ ticks, soit environ 5 cents, c'est-à-dire
$4.6/8\,000 = \qty{5.8}{\bp}$.
\end{solution}

\begin{solution}{exo:m1:auctions:6}
À la clôture~: $0.12\% \times 50 = \$60\,000$ au-dessus du cours actuel. Maintenant~:
$0.08\% \times 50 = \$40\,000$ d'impact, \emph{auquel s'ajoute} le risque que la clôture
diffère du prix obtenu, un risque que le fonds ne saurait justifier~: son
indice de référence est la clôture, et un achat à la clôture a une \omterm{def:m1:the-buy-side:benchmark}{tracking error} nulle
quelle que soit la prime. Ces 20\,000~\$ représentent le prix de cette certitude, payé par
les investisseurs du fonds, et invisible dans ses statistiques de suivi.
\end{solution}

\begin{solution}{exo:m1:auctions:7}
Avec 400 ordres, un déséquilibre de 50~\% (environ 45\,000 actions) déplace la clôture
de 1.83 ticks~: 4.1 ticks pour 100\,000 actions. Avec 1\,600 ordres, le même
50~\% correspond à environ 182\,000 actions et la déplace de 1.87 ticks~: 1.0 tick pour 100\,000
actions. Quatre fois plus d'ordres signifie quatre fois plus de densité d'offre et de
demande par tick, d'où un quart du déplacement par action, comme le
prédit la \cref{prop:m1:auctions:depth}~; un déséquilibre qui représente une
\emph{fraction} fixe de l'\omterm{def:m1:auctions:call}{enchère} déplace le prix du même montant dans les deux cas.
\end{solution}

\begin{solution}{exo:m1:auctions:8}
(i) Au \omterm{def:m1:auctions:call}{prix d'équilibre}, les ordres sont exécutés selon la priorité temporelle, et le
côté en excédent n'est exécuté qu'en partie~: un ordre saisi à 15:55 exactement au
\omterm{def:m1:auctions:moc}{cours de clôture} se place derrière tous les ordres antérieurs et reçoit souvent rien. (ii) Le
\omterm{def:m1:auctions:imbalance}{cours indicatif} à 15:55 est lui-même déplacé par l'ordre de la stratégie et par les
intérêts de contrepartie qu'il attire~; lorsque la clôture passe \emph{au-delà}
de la limite, c'est en général parce qu'un intérêt mieux informé est arrivé ensuite,
si bien que les exécutions obtenues subissent une \omterm{def:m1:what-a-trading-firm-does:adverse-selection}{sélection adverse}. Corriger cela suppose de connaître les
volumes appariés et le déséquilibre de l'\omterm{def:m1:auctions:call}{enchère} au \omterm{def:m1:auctions:call}{prix d'équilibre}, la position de l'ordre
dans la file temporelle à sa limite, et une estimation de l'impact propre de la
stratégie, à partir de la \cref{prop:m1:auctions:depth}.
\end{solution}

\begin{solution}{pb:m1:auctions:1}
\textbf{1.} $D$~: 400\,000 à 39.90~; 250\,000 à 39.95 et 40.00~; 130\,000 à
40.05 et 40.10~; 50\,000 à partir de 40.15. $S$~: 60\,000~; 160\,000~; 160\,000~;
300\,000~; 300\,000~; 500\,000 à 40.15~; 750\,000 à 40.30~; 1\,050\,000 à
40.50.
\textbf{2.} Volume exécutable~: 60, 160, 160, 130, 130, 50, 50, 50 (milliers). Le
maximum, 160\,000, est atteint à 39.95 et 40.00 avec le même excédent~; dans les deux cas,
l'excédent est à l'achat, donc la règle 3 retient le plus élevé~: 40.00.
\textbf{3.} $250\,000 - 160\,000 = 90\,000$ à l'achat.
\textbf{4.} L'ordre d'achat à 40.00 reçoit 30\,000 de ses 120\,000.
\textbf{5.} Chaque valeur de $D$ augmente de 500\,000~: 900, 750, 750, 630, 630,
550, 550, 550 (milliers).
\textbf{6.} Volume exécutable~: 60, 160, 160, 300, 300, 500, 550, 550. Le maximum
550\,000 est atteint à 40.30 et 40.50~; l'excédent est de 200\,000 à vendre à 40.30
et 500\,000 à 40.50~: la règle 2 donne 40.30.
\textbf{7.} 30 cents, 75 points de base.
\textbf{8.} À 40.00~: $750\,000 - 160\,000 = 590\,000$ actions à acheter.
\textbf{9.} $0.30 \times 500\,000 = \$150\,000$.
\textbf{10.} $S$ passe à 450\,000 à 40.10, 650\,000 à 40.15, 850\,000 à
40.20, 1\,100\,000 à 40.30. Volume exécutable de 550\,000 à partir de 40.15~; le
plus petit excédent, 100\,000 à vendre, se trouve à 40.15~: la clôture est à 40.15 pour
550\,000 actions.
\textbf{11.} $0.15 \times 500\,000 = \$75\,000$.
\textbf{12.} Il a vendu à 40.15 et compte racheter à 40.05~: $0.10 \times
150\,000 = \$15\,000$.
\textbf{13.} Il est vendeur à découvert de 150\,000 actions pour la nuit, sur un titre qui vient
de faire son entrée dans un indice~: 6~millions de dollars de risque sur un seul nom, y compris toute annonce. Il peut
couvrir la composante de marché avec des contrats à terme sur indice, ou avoir acheté les actions
au préalable, pendant les jours séparant l'annonce de l'inclusion --- auquel
cas sa vente à la clôture complète une transaction d'inventaire.
\textbf{14.} Entre 40.00 et 40.15, l'offre augmente de 490\,000 et la demande
baisse de 200\,000~: environ 46\,000 actions par tick. La proposition prédit
$500\,000/46\,000 = 10.9$ ticks, contre 15 observés~: le carnet est irrégulier,
et l'ordre lié à l'indice a épuisé la zone dense voisine de 40.00.
\textbf{15.} Un prix moyen plus bas, en anticipant leur propre demande~; au
prix d'une \omterm{def:m1:the-buy-side:benchmark}{tracking error} par rapport à un indice qui n'intègre le titre qu'à partir de
la clôture de ce soir, et de révéler leurs achats à ceux qui négocient avant
eux, de toute façon.
\textbf{16.} Les investisseurs des fonds indiciels, à travers un \omterm{def:m1:auctions:moc}{cours de clôture} supérieur de 15 cents
au-dessus du marché d'avant l'\omterm{def:m1:auctions:call}{enchère}~: les 15\,000~\$ des fournisseurs de liquidité en font partie, les
vendeurs des ordres à cours limité exécutés pendant la hausse perçoivent le reste.
\textbf{17.} Pour faire monter le \omterm{def:m1:auctions:imbalance}{cours indicatif} et inciter d'autres à présenter des ordres de vente ou à
renoncer à enchérir, puis annuler ses ordres~: il n'a jamais eu l'intention de négocier.
Il s'agit de spoofing, une forme de manipulation de marché~; à la clôture, il
fausse aussi le \omterm{def:m1:auctions:moc}{cours de clôture}.
\textbf{18.} Afin que le déséquilibre publié soit une grandeur crédible qui ne puisse
que diminuer~: des ordres tardifs susceptibles de l'accroître, ou des annulations
pourraient le faire réapparaître, ce qui rendrait l'information sans valeur et le \omterm{def:m1:auctions:moc}{cours de clôture}
facile à manipuler.
\textbf{19.} \textbf{40.15 pour 550\,000 actions.}
\textbf{20.} La certitude de détenir le titre exactement au prix que
utilise son indice de référence.
\end{solution}

\begin{solution}{iq:m1:auctions:1}
Les ordres s'accumulent sans donner lieu à transaction. À la clôture, la bourse détermine le prix
qui maximise le volume apparié~; les égalités sont départagées par le plus petit déséquilibre,
puis par le sens du déséquilibre, puis par la proximité avec un prix de référence.
Tous les ordres appariés sont exécutés à ce prix unique, les ordres au marché et les ordres à cours limité
les plus favorables passant en premier, puis la priorité temporelle au \omterm{def:m1:auctions:call}{prix d'équilibre}.
\iqlookfor{volume maximal d'abord, et une séquence raisonnable de départage des égalités.}
\end{solution}

\begin{solution}{iq:m1:auctions:2}
Les fonds indiciels et les gérants évalués par rapport à un indice valorisent au \omterm{def:m1:auctions:moc}{cours de clôture} et évitent
le \omterm{def:m1:the-buy-side:benchmark}{tracking error} en y négociant~; les produits dérivés et les flux des fonds s'y dénouent~; la
\omterm{def:m1:auctions:call}{enchère} offre la liquidité la plus profonde de la journée, sans payer de fourchette~; et le volume qui
s'y déplace amincit le marché intrajournalier et attire encore plus de volume vers ce moment. La
\omterm{def:m1:the-buy-side:passive}{gestion passive} renforce chacun de ces effets chaque année.
\iqlookfor{le recours aux indices de référence comme cause première.}
\end{solution}

\begin{solution}{iq:m1:auctions:3}
Trier les ordres d'achat à cours limité par ordre décroissant et les ordres de vente à cours limité par ordre croissant, les ordres au marché
en tête de chaque côté~; parcourir une seule fois les prix distincts avec les quantités demandées
et offertes cumulées (deux pointeurs), en retenant le volume exécutable maximal et, parmi les égalités,
le surplus minimal~: $O(n\log n)$ pour les tris, $O(n)$ pour le
balayage. Cas limites~: pas de croisement~; seuls des ordres au marché des deux côtés (aucun
prix limite~: utiliser la référence)~; égalités sur une plage de prix~; plages à surplus nul~;
quantités qui dépassent 32 bits~; allocation au \omterm{def:m1:auctions:call}{prix d'équilibre} lorsque le côté
en surplus doit être réparti au prorata ou selon la priorité temporelle~; prévention de l'auto-appariement.
\iqlookfor{le balayage unique avec sommes cumulées, et le cas où seuls les ordres au marché
sont présents.}
\end{solution}

\begin{solution}{iq:m1:auctions:4}
Le déséquilibre dépasse une \omterm{def:m1:auctions:call}{demi-enchère} moyenne~: le \omterm{def:m1:auctions:imbalance}{cours indicatif} sera nettement au-dessus
de la dernière transaction. Des intérêts vendeurs compensateurs arrivent ---
fournisseurs de liquidité et investisseurs ayant passé des \omterm{def:m1:auctions:moc}{ordres à cours limité à la clôture} (LOC) ---
réduisant le déséquilibre et ramenant le \omterm{def:m1:auctions:imbalance}{cours indicatif} vers le marché, mais
pas complètement. Le cours continu est tiré à la hausse dans les dernières minutes
comme ceux qui vendent à l'\omterm{def:m1:auctions:call}{enchère} se couvrent en achetant dès maintenant, et comme d'autres achètent
par anticipation avant la clôture pour revendre dans l'\omterm{def:m1:auctions:call}{enchère}. On peut s'attendre à un retournement partiel à l'ouverture suivante.
\iqlookfor{à la fois la dynamique de l'\omterm{def:m1:auctions:call}{enchère} et les effets de débordement sur le
carnet continu.}
\end{solution}

\begin{solution}{iq:m1:auctions:5}
Cible~: \omterm{def:m1:auctions:moc}{cours de clôture} moins milieu de fourchette au moment de la prévision (en points de base ou
en fourchettes), ou retournement du lendemain matin. Variables~: déséquilibre signé rapporté
au volume moyen des \omterm{def:m1:auctions:call}{enchères} et à la profondeur affichée, quantité appariée, \omterm{def:m1:auctions:imbalance}{cours indicatif}
rapporté au milieu, variations entre messages, temps avant la clôture,
événements d'indice, volatilité. Piège~: le flux de déséquilibre et les cotations continues
arrivent avec des latences différentes, et les variables construites à partir de messages publiés
\emph{après} l'instant de la décision révèlent la réponse~; l'évaluation doit aussi tenir compte
de l'impact de l'ordre de la stratégie lui-même sur le déséquilibre qu'elle conditionne.
\iqlookfor{biais d'anticipation via les horodatages, et impact propre.}
\end{solution}

\begin{solution}{iq:m1:auctions:6}
Gagnants~: les fournisseurs de liquidité plus lents, qui ne peuvent plus subir de sniping pendant un
intervalle de traitement par lots, puisque tous les ordres arrivant dans les mêmes 100 millisecondes
se disputent le prix et non le temps~; les investisseurs, si l'économie est répercutée
sous forme de fourchettes plus étroites. Perdants~: les acteurs dont l'avantage est la seule vitesse~; éventuellement la plateforme,
si les preneurs de liquidité n'aiment pas attendre et partent vers des concurrents en continu. Les fourchettes
devraient se resserrer là où le coût du sniping représentait une grande part d'entre elles, et ne s'élargir nulle part~;
mais, si une seule plateforme l'adopte, l'arbitrage entre elle et les plateformes continues
réintroduit une course à chaque frontière de lot, et le bénéfice dépend de la part
du marché qui se déplace.
\iqlookfor{la concurrence par les prix remplaçant la concurrence par le temps, et la
mise en garde sur la fragmentation.}
\end{solution}
